% !TEX program = xelatex
% ============================================================
%  Arabic Work Breakdown Structure (WBS) Template — Nibras Center
%  مشروع بحثي: هيكل تقسيم العمل
%  Compile with: xelatex main.tex (twice)
% ============================================================

\documentclass[11pt,a4paper]{article}

\usepackage[a4paper,margin=2cm,top=2.5cm,bottom=2.5cm,headheight=15pt]{geometry}
\usepackage{amsmath}
\usepackage{graphicx}
\usepackage{array}
\usepackage{booktabs}
\usepackage{tabularx}
\usepackage{multirow}
\usepackage{enumitem}
\usepackage{float}
\usepackage{titlesec}
\usepackage{fancyhdr}
\usepackage{setspace}
\usepackage[table]{xcolor}
\usepackage{tikz}
\usetikzlibrary{positioning,calc,shapes.geometric,arrows.meta}

\usepackage{bidi}
\usepackage[no-math]{fontspec}
\usepackage[quiet]{polyglossia}

\setmainfont[Scale=1]{NotoKufiArabic-Regular.ttf}
\newfontfamily\arabicfont[Script=Arabic,Scale=0.95]{NotoKufiArabic-Regular.ttf}
\newfontfamily\englishfont[Scale=0.95]{TeX Gyre Termes}

\setdefaultlanguage[calendar=gregorian,hijricorrection=1]{arabic}
\setotherlanguage[variant=british]{english}

\usepackage[breaklinks,hidelinks]{hyperref}
\hypersetup{pdftitle={هيكل تقسيم العمل},pdfauthor={الباحث الرئيسي}}

\titleformat{\section}{\Large\bfseries}{\thesection}{0.7em}{}
\titlespacing*{\section}{0pt}{1.4ex plus 0.4ex}{0.9ex plus 0.3ex}

\pagestyle{fancy}
\fancyhf{}
\fancyhead[R]{\small\itshape هيكل تقسيم العمل (WBS)}
\fancyhead[L]{\small اسم المشروع}
\fancyfoot[C]{\small\thepage}
\renewcommand{\headrulewidth}{0.3pt}

\newcolumntype{L}[1]{>{\raggedright\arraybackslash}p{#1}}
\newcolumntype{C}[1]{>{\centering\arraybackslash}p{#1}}

\definecolor{wbsLevel1}{HTML}{1A3C6B}
\definecolor{wbsLevel2}{HTML}{2E6DA4}
\definecolor{wbsLevel3}{HTML}{D4A017}
\definecolor{wbsBox}{HTML}{F0F4F8}

\setstretch{1.3}

\begin{document}

\begin{center}
{\LARGE\bfseries هيكل تقسيم العمل (WBS)}\\[0.5em]
{\Large عنوان المشروع البحثي}\\[1em]
\end{center}

% ============================================================
% 1. WBS TABLE
% ============================================================
\section{جدول هيكل تقسيم العمل}
\label{sec:wbs_table}

\begin{table}[H]
\centering
\small
\renewcommand{\arraystretch}{1.3}
\begin{tabularx}{\textwidth}{C{1.5cm} L{5cm} L{3cm} C{2cm} X}
\toprule
\textbf{الرمز} & \textbf{اسم الحزمة/المهمة} & \textbf{المسؤول} & \textbf{المدة} & \textbf{المخرجات} \\
\midrule
\rowcolor{wbsBox}
\textbf{1.0} & \textbf{المرحلة الأولى: المراجعة والتخطيط} & الباحث الرئيسي & 3 أشهر & تقرير الأدبيات \\
1.1 & مراجعة الأدبيات السابقة & باحث مشارك & 6 أسابيع & فصل الأدبيات \\
1.2 & تحديد الفجوات البحثية & الباحث الرئيسي & 2 أسبوع & وثيقة الفجوات \\
1.3 & تصميم المنهجية & الباحث الرئيسي & 4 أسابيع & وثيقة المنهجية \\
\midrule
\rowcolor{wbsBox}
\textbf{2.0} & \textbf{المرحلة الثانية: الإعداد والتجربة} & باحث مشارك & 4 أشهر & البيانات الخام \\
2.1 & تجهيز المعدات والأدوات & فني المختبر & 3 أسابيع & قائمة التجهيز \\
2.2 & جمع البيانات & باحث مشارك & 10 أسابيع & قاعدة البيانات \\
2.3 & معالجة البيانات & باحث مشارك & 3 أسابيع & بيانات منظفة \\
\midrule
\rowcolor{wbsBox}
\textbf{3.0} & \textbf{المرحلة الثالثة: التحليل والنمذجة} & الباحث الرئيسي & 4 أشهر & نتائج التحليل \\
3.1 & التحليل الإحصائي & الباحث الرئيسي & 6 أسابيع & جداول النتائج \\
3.2 & النمذجة والمحاكاة & باحث مشارك & 6 أسابيع & نماذج \\
3.3 & التحقق من النتائج & الباحث الرئيسي & 4 أسابيع & تقرير التحقق \\
\midrule
\rowcolor{wbsBox}
\textbf{4.0} & \textbf{المرحلة الرابعة: النشر والتوثيق} & الباحث الرئيسي & 5 أشهر & الورقة والتقرير \\
4.1 & كتابة الورقة البحثية & الباحث الرئيسي & 8 أسابيع & مسودة الورقة \\
4.2 & تقديم الورقة للنشر & الباحث الرئيسي & 4 أسابيع & تأكيد النشر \\
4.3 & كتابة التقرير النهائي & الباحث الرئيسي & 6 أسابيع & التقرير النهائي \\
4.4 & عرض النتائج & الفريق كله & 2 أسبوع & عرض تقديمي \\
\bottomrule
\end{tabularx}
\end{table}

% ============================================================
% 2. WBS DIAGRAM (Tree)
% ============================================================
\section{المخطط الشجري لهيكل تقسيم العمل}
\label{sec:wbs_diagram}

\begin{figure}[H]
\centering
\begin{tikzpicture}[
  every node/.style={font=\small},
  level1/.style={fill=wbsLevel1,text=white,rounded corners=3pt,inner sep=6pt,minimum width=3.5cm,align=center},
  level2/.style={fill=wbsLevel2,text=white,rounded corners=3pt,inner sep=5pt,minimum width=2.8cm,align=center},
  level3/.style={fill=wbsBox,draw=wbsLevel2,thick,rounded corners=3pt,inner sep=4pt,minimum width=2.2cm,align=center},
  line/.style={draw=wbsLevel2,thick},
]

% Root
\node[level1] (root) at (0,0) {المشروع البحثي};

% Level 1 — 4 phases
\node[level2] (p1) at (-7.5,-2.5) {1. المراجعة والتخطيط};
\node[level2] (p2) at (-2.5,-2.5) {2. الإعداد والتجربة};
\node[level2] (p3) at (2.5,-2.5) {3. التحليل والنمذجة};
\node[level2] (p4) at (7.5,-2.5) {4. النشر والتوثيق};

% Lines from root to phases
\draw[line] (root) -- (p1);
\draw[line] (root) -- (p2);
\draw[line] (root) -- (p3);
\draw[line] (root) -- (p4);

% Level 2 — sub-tasks
\node[level3] (p11) at (-9.5,-5) {1.1 مراجعة الأدبيات};
\node[level3] (p12) at (-6.5,-5) {1.2 الفجوات البحثية};
\node[level3] (p13) at (-3.5,-5) {1.3 تصميم المنهجية};

\node[level3] (p21) at (-5.5,-6.5) {2.1 التجهيز};
\node[level3] (p22) at (-2.5,-6.5) {2.2 جمع البيانات};
\node[level3] (p23) at (0.5,-6.5) {2.3 المعالجة};

\node[level3] (p31) at (0.5,-5) {3.1 التحليل الإحصائي};
\node[level3] (p32) at (3.5,-5) {3.2 النمذجة};
\node[level3] (p33) at (6.5,-5) {3.3 التحقق};

\node[level3] (p41) at (5.5,-6.5) {4.1 كتابة الورقة};
\node[level3] (p42) at (8,-6.5) {4.2 التقديم للنشر};
\node[level3] (p43) at (10.5,-6.5) {4.3 التقرير};

\draw[line] (p1) -- (p11);
\draw[line] (p1) -- (p12);
\draw[line] (p1) -- (p13);

\draw[line] (p2) -- (p21);
\draw[line] (p2) -- (p22);
\draw[line] (p2) -- (p23);

\draw[line] (p3) -- (p31);
\draw[line] (p3) -- (p32);
\draw[line] (p3) -- (p33);

\draw[line] (p4) -- (p41);
\draw[line] (p4) -- (p42);
\draw[line] (p4) -- (p43);

\end{tikzpicture}
\caption{المخطط الشجري لهيكل تقسيم العمل}
\end{figure}

% ============================================================
% 3. WORK PACKAGES DETAIL
% ============================================================
\section{تفاصيل حزم العمل}
\label{sec:work_packages}

\subsection{حزمة العمل 1.0: المراجعة والتخطيط}
\textbf{المسؤول:} الباحث الرئيسي \quad \textbf{المدة:} 3 أشهر\\
\textbf{المخرجات:} تقرير الأدبيات، وثيقة الفجوات، وثيقة المنهجية\\
\textbf{معايير القبول:} مراجعة محكمة من قبل المشرف

\subsection{حزمة العمل 2.0: الإعداد والتجربة}
\textbf{المسؤول:} باحث مشارك \quad \textbf{المدة:} 4 أشهر\\
\textbf{المخرجات:} قاعدة بيانات منظفة\\
\textbf{معايير القبول:} اكتمال جميع العينات وجودة البيانات

\subsection{حزمة العمل 3.0: التحليل والنمذجة}
\textbf{المسؤول:} الباحث الرئيسي \quad \textbf{المدة:} 4 أشهر\\
\textbf{المخرجات:} جداول النتائج والنماذج\\
\textbf{معايير القبول:} دقة إحصائية مقبولة وتحقق مستقل

\subsection{حزمة العمل 4.0: النشر والتوثيق}
\textbf{المسؤول:} الباحث الرئيسي \quad \textbf{المدة:} 5 أشهر\\
\textbf{المخرجات:} ورقة منشورة وتقرير نهائي\\
\textbf{معايير القبول:} قبول الورقة للنشر وموافقة الجهة الممولة

\end{document}
